\section{Scope and Assumptions}
\todo{Trim down this sections. It's way too verbose}
This thesis is concerned with the \emph{software design and evaluation} of radar tracking pipelines for UAV detect-and-avoid under size, weight and power constraints. Its scope is limited to the processing stages that occur after radar samples have been acquired, with emphasis on how alternative signal-processing, angle-estimation, and tracking algorithms perform when implemented on embedded compute platforms suitable for small and medium UAVs. The work therefore focuses on hardware-aware software design rather than radar hardware design itself. In particular, the thesis evaluates candidate pipelines through three main forms of analysis: profiling of individual algorithms, profiling of integrated end-to-end pipelines, and sensitivity analysis of detection and tracking performance under constrained signal conditions. These evaluations are used to compare cost and performance in terms of latency, throughput, memory behaviour, and sensitivity to faint or distant targets. 

Within this scope, the processing chain considered includes front-end shaping and windowing, fast Fourier transforms, coherent integration, constant false alarm rate detection, monopulse angle sensing, coordinate transformation, data association, tracking filters, and track management. The thesis also considers orchestration and dataflow decisions, such as batching, reading and writing strategies, stream aggregation, and the performance implications of host--device transfers and synchronisation. These elements are included because the identified knowledge gap is not only algorithmic, but systems-level: the suitability of a radar tracking pipeline for UAV detect-and-avoid depends jointly on sensing performance and implementation cost on embedded hardware. 

\todo{This is a bit repetative}
The empirical scope includes both controlled and partially realistic evaluation. Candidate methods are tested using a combination of simulation, synthetic or simplified input data, and captured radar data where available, in order to study both low-level computational behaviour and tracking performance in encounter-relevant scenarios. Partial field testing is included only to the extent necessary to demonstrate practical feasibility and to expose the pipeline to non-ideal data. The thesis does not attempt a complete flight-qualified validation campaign, nor does it claim certification readiness for operational UAV detect-and-avoid deployment. Rather, the aim is to provide comparative evidence and design guidance for practical onboard radar processing under realistic embedded constraints. 

\todo{Idk if I need the word "explicitly"}
Several topics are explicitly outside the scope of this work. First, the thesis does not address the design of the radar front-end hardware, including antenna design, RF electronics, waveform generation, or board-level integration. Second, it does not provide full-scale field trials across all weather, geometry, clutter, and airspace conditions that would be required for operational certification. Third, although FPGA- and ASIC-based implementations are relevant to radar systems, they are not the focus here; the work is restricted to programmable CPU/GPU-oriented platforms in order to study software-defined architectures that can be iterated rapidly and deployed on modern embedded UAV compute stacks. Finally, this thesis does not attempt a cross-modality comparison between radar and optical or infrared sensing. Optical and infrared systems are discussed only to motivate radar's complementary role, not to benchmark relative sensing performance across modalities. 

The thesis proceeds under several assumptions. The first is that the radar sensor is a \emph{small conformal monopulse radar} whose aperture and power limitations are representative of a UAV-compatible installation. This assumption is consistent with the thesis focus on low-SNR operation and reduced angular resolution relative to larger airborne radar systems. The second is that the target compute platform includes \emph{GPU-enabled embedded hardware}, reflecting the increasing use of heterogeneous CPU/GPU companion computers in modern UAV systems and the thesis emphasis on highly parallel radar processing. The third is that the hardware context of primary interest includes \emph{shared or unified memory} between CPU and GPU, as found in many embedded systems-on-module, although comparison against CPU-only and discrete-memory contexts may be used where available to clarify architectural trade-offs. These assumptions are adopted because they reflect the practical hardware class for which the pipeline is being designed: not a large airborne radar platform with abundant power and aperture, but a compact UAV-compatible software-defined radar system in which memory movement, synchronisation overhead, and implementation structure materially affect performance. 

Under these boundaries, the thesis does not seek to produce a universally optimal radar architecture. Instead, it seeks to identify which combinations of algorithms and implementation strategies are most suitable for the stated UAV context, and how those conclusions vary with embedded hardware restrictions. The outcome is intended to be a set of hardware-aware recommendations for the design of practical radar tracking pipelines for UAV detect-and-avoid, rather than a complete end-to-end certified system design. 